% Getting Started, A Brief Lesson (AN-21 F-WPC; model: WordPerfect 5.1
% Reference pp. 18-21: type, save, exit).  Every keystroke below is
% replayed on WORD.EXE by `make test-pubs-iword` and every screen panel is
% the screen that replay produces (GC-03; initech-wpc.sty \BeginReplay).
% GATES: test-pubs-iword test-iword-face
\Gates{test-pubs-iword test-iword-face}
\WPHead{A Brief Lesson}
In this lesson you type a short letter, make one line bold, save the letter
and leave InitechWord.  Start InitechWord as described under
\textit{Starting InitechWord} above, so that you have a new document on the
screen.

\BeginReplay{gs-brief}
\WPSide{Typing the Letter}
\begin{WPSteps}
\item Type the salutation:
\Type{Dear Mr. Lumbergh:}
\item Press \PressKey{Enter} twice to end the line and leave a blank line.\Keys{Enter}
\item Type the paragraph below.  Do not press Enter at the end of a line:
InitechWord moves to the next line when the line is full.
\Type{Here is the summary you asked for. A cover sheet has gone out with every TPS report this week.}
\item Press \PressKey{Enter} twice.\Keys{Enter}
\end{WPSteps}

\WPSide{Making a Line Bold}
\begin{WPSteps}
\item Press \Press{Bold}{F6}.  Nothing seems to change, but InitechWord has
put a pair of bold codes in the document, with the cursor between them.
\item Type the closing:
\Type{Thank you.}
\item Press \Press{Bold}{F6} again.  The cursor moves past the code that
ends the bold, so that what you type next is not bold.
\end{WPSteps}
Your screen now looks like this.  The closing is shown in bold, and the
status line shows that the cursor is on the line 1.83 inches from the top
of the page, two inches from the left edge.
\WordScreen[rows=6]{gs-brief-1}{The letter, with the closing in bold}

\WPSide{Saving the Letter}
\begin{WPSteps}
\item Press \Press{Save}{F10}.  The status line asks for a name.
\PullDown{Select \Mi{S}ave from the \Mi{F}ile menu.}
\item Type \TypeName{letter} as the name and press \PressKey{Enter}.
\end{WPSteps}
InitechWord saves the letter on the data volume and shows its name at the
left of the status line.  Names are always shown in capitals.
\WordScreen[rows=6]{gs-brief-2}{The saved letter}

\WPSide{Leaving InitechWord}
\begin{WPSteps}
\item Press \Press{Exit}{F7}.  InitechWord asks whether to save the
document.  At the right of the status line it adds that the text has not
been changed since you saved it.
\WordScreen[rows=6]{gs-brief-3}{Exit asks whether to save the document}
\item Type \TypeName{n}.  The letter is already saved.  InitechWord asks
whether you want to leave the program.
\item Type \TypeName{y} to leave InitechWord.
\end{WPSteps}
\EndReplay
The InitechWord window closes and you are back at the FLAIR desktop.  The
letter is in the file \Scr{LETTER} on the data volume.  The
\textit{Workbook} lessons continue from here.
